\documentclass{article}

\usepackage[english]{babel}
\usepackage[utf8x]{inputenc}
\usepackage{amsmath}
\usepackage{amsthm}
\usepackage{graphicx}
\usepackage[colorinlistoftodos]{todonotes}

\title{Initial \LaTeX exercise, CSci 1302}
\author{\emph{Replace this with your name}}

\begin{document}
\maketitle

\todo[inline]{Make sure to put your name up in the "Author" tag above.}

\todo[inline]{When you've finished the problem set you should remove all the comments/hints (the big orange boxes of text) that I've provided using the \texttt{\textbackslash todo} tag.}

\section{Welcome to \LaTeX!}

"Fill in the blanks" below to answer a few questions for me, and as a way for us all to get started with \LaTeX and make sure this works. In \LaTeX's simplest form, all you have to do is write text, using blank lines to separate paragraphs. 

\begin{center}
\em{Thanks for your time!} -- Nic
\end{center}

\section{How would you prefer to be addressed?}f

\todo[inline]{Replace the paragraph below with a short paragraph telling me your name and preferred pronouns. You don't have to say anything about how your name is pronounced, but you certainly can if you wish. I'd love it if you e-mailed me a little voice memo demonstrating how you pronounce your name, and that can be much more useful than a written description of the pronunciation.}

The name I prefer is "Nic" (pronounced like "Nick", but written without the 'k' for murky reasons having to do with Swiss-German), and my preferred pronouns are "he/him/his".

\section{Why are you taking CSci 1302?}

\todo[inline]{Write a sentence or three about why you're taking this class and what you hope to get out of it. It's totally OK to just say you're taking it because it's required for the CSci major, but it would be nice if you have a look at the syllabus and think a little about \em{why} we require it and what it might do for you.}

\section{Any questions?}

\todo[inline]{Use this section to ask any questions you might have about the course, about the CSci major/minor, about computing in general, etc., etc.}

\end{document}